\section{A2A tasks artifacts and delegated autonomy}\label{sec:C19:a2a}
A2A allows a remote agent to respond with a self-contained message or initiate stateful work represented as a task. Its lifecycle distinguishes interrupted states, such as \texttt{input-required} or \texttt{auth-required}, from terminal outcomes. A context identifier can group multiple interactions, while a task identifier refers to a particular unit of work~\cite{a2a2026task}. These identifiers should not be mapped automatically to the model's active context window or to one transport connection.

Terminal tasks do not restart. A refinement creates new work within the related context and can reference the earlier task~\cite{a2a2026task}. Consequently, client artifact acceptance and mutation lineage must survive beyond the producing task's terminal status. Reusing the old task identifier as a mutable revision slot would obscure which input produced which artifact. Continuity belongs to the explicit context and client-managed lineage, while terminal task identity preserves the completed unit of work.

A2A v1.0.1 also distinguishes idempotency by operation. Reads are naturally idempotent, message submission MAY use the message identifier for duplicate detection, and task cancellation is idempotent~\cite{cumb2026a2aspec101}. Thus a repeated message identifier is not a universal delivery guarantee, and idempotent cancellation does not establish rollback of a previously committed effect. The client must record which remote deduplication contract was actually promised.

Discovery uses an Agent Card describing offered capabilities, endpoints, input and output modes, and security-related information~\cite{cumb2026a2adiscovery101}. The card supports selection of a suitable remote agent; it does not prove the agent's competence. A skill description in the card is an advertised capability, not necessarily a portable executable skill package. The caller still needs task validation, evidence handling, and a policy for trusting returned artifacts.

Delegating a task differs from invoking a narrow tool because the remote agent can choose internal methods and produce intermediate work. The caller may observe only status and artifacts rather than every internal action. This can be useful for modularity and privacy, but it changes verification. A completed task state is a protocol statement from the remote service, not independent proof that the original user's success predicate was met.

Artifacts should preserve provenance, version, and a relation to the task. A document returned by a remote researcher can contain supported claims, unsupported synthesis, or stale citations. An execution artifact can be syntactically valid while violating a local constraint unknown to the remote agent. The delegation request should therefore include the relevant obligations and authority, and the caller should validate the artifact against them before publication or further effects.

Parallel follow-ups create a state-ownership problem. Several tasks can share a context identifier while working on different branches of an objective. A late artifact from one branch should not silently overwrite a newer accepted artifact from another. The application needs version and acceptance rules beyond the protocol's ability to transport updates. Similarly, a request for more information should suspend or qualify dependent work rather than being ignored while the caller assumes progress.

A proposed delegation contract contains goal, completion predicate, budget, permitted effects, evidence requirements, and expected artifact schema. It allows the remote agent flexibility within those constraints. A callback or streaming update can report progress, while final acceptance remains a separate local decision. This distinguishes negotiated autonomy from unrestricted trust and connects protocol design to the orchestration analysis in Chapter~\ref{ch:C21}.

